\addcontentsline{toc}{chapter}{Resumen}

\chapter*{Resumen}

\vspace*{-8mm}

Se presenta una implementación de Redes Neuronales Informadas por Física con
activación sinusoidal (PINN-SIREN) para la validación de la resolución de la
ecuación de Helmholtz en una y dos dimensiones, como base para la simulación
de speckle óptico.
Con $\omega_0 = 1.0$ calibrado al número de onda, la arquitectura SIREN
alcanza un error $L^2$ de $0.006\%$ en el caso 1D ($5 \times 64$ neuronas) y
de $0.171\%$ en el caso 2D ($5 \times 128$) respecto a soluciones analíticas,
tres y casi dos órdenes de magnitud por debajo del umbral de tesis ($5\%$).
La extensión a speckle se formula mediante una pantalla de fase aleatoria
espacialmente correlacionada y una referencia semianalítica de propagación
por espectro angular, y se resuelve con una reducción modal que impone la
condición de Cauchy de forma exacta en la arquitectura. En cinco
realizaciones de fase independientes a $z=1\lambda$ el error $L^2$ del campo
complejo es de $3.161\%$, dominado por la fracción evanescente del espectro
que la base modal no puede representar; medido sobre el subespacio
propagante, el error es de $0.042\%$. El contraste del conjunto es
$C=0.9489$. El contraste de referencia promedio de los conjuntos de 64
realizaciones es $C=0.9661$.
Una estrategia de optimización bifásica Adam~+~L-BFGS y el muestreo por
hipercubo latino (LHS) se emplean en el protocolo de entrenamiento.
La reproducibilidad del método fue verificada con tres semillas independientes,
obteniendo un promedio componente a componente de error $L^2$ de
$0.192 \pm 0.089\%$ (desviación estándar poblacional).

\bigskip

\noindent\textbf{Palabras clave:} Redes Neuronales Informadas por Física,
SIREN, ecuación de Helmholtz, speckle óptico, activación sinusoidal,
muestreo por hipercubo latino.
